% =============================================================================
%  Per-result proof paper -- TEMPLATE
%  Copy this directory to docs/papers/<id>_<slug>/, then fill in:
%    - \title, \author
%    - the Statement (verbatim from the scaffolding paper)
%    - the Proof
%    - the Discussion
%  Build with: just paper-result <id>_<slug>
% =============================================================================
\documentclass[11pt]{article}

% ---------- Encoding & fonts ----------
\usepackage[utf8]{inputenc}
\usepackage[T1]{fontenc}
\usepackage[expansion=false]{microtype}

% ---------- Math ----------
\usepackage{amsmath, amssymb, amsthm, mathtools}

% ---------- Layout ----------
\usepackage[margin=1in]{geometry}
\usepackage{enumitem}

% ---------- Hyperlinks ----------
\usepackage[hidelinks]{hyperref}

% ---------- Theorem environments ----------
\theoremstyle{plain}
\newtheorem{theorem}{Theorem}
\newtheorem{proposition}[theorem]{Proposition}
\newtheorem{lemma}[theorem]{Lemma}
\newtheorem{corollary}[theorem]{Corollary}

\theoremstyle{definition}
\newtheorem{definition}[theorem]{Definition}

\theoremstyle{remark}
\newtheorem*{remark}{Remark}

% =============================================================================
\title{TEMPLATE: Result on \texttt{<conjecture-or-open-problem-id>}}
\author{Author Name\thanks{Affiliation. Email.}}
\date{\today}

\begin{document}
\maketitle

\begin{abstract}
One paragraph: which conjecture or open problem this paper addresses, what is
established (proved, disproved, partially resolved, reformulated), and the
main idea of the argument.
\end{abstract}

% -----------------------------------------------------------------------------
\section{Statement}

Reproduce the conjecture or open problem \emph{verbatim} from the scaffolding
paper~\cite{emergent_systems_2026}, with a citation to the relevant section.

\begin{quote}
\emph{[Verbatim statement here, e.g.\ ``C1 (closure-operator unification of
viability): each of the four viability formalisms\dots''.]}
\end{quote}

\noindent (Scaffolding paper, \S5.1, statement~\texttt{<id>}.)

% -----------------------------------------------------------------------------
\section{Result}

State the result of this paper as a single proposition or theorem, in the
notation of the scaffolding paper. If the result is partial, say so explicitly
and characterise the gap.

\begin{theorem}[Main result]
\textit{Statement of the result.}
\end{theorem}

% -----------------------------------------------------------------------------
\section{Proof}

\begin{proof}
The proof.
\end{proof}

% -----------------------------------------------------------------------------
\section{Discussion}

What does this result mean for the framework? Does it close the conjecture, or
does it reformulate it (in which case, link the new statement and update
\texttt{RESULTS.md} with status \texttt{mutated}). What follow-up questions
does it open?

% =============================================================================
\bibliographystyle{plain}
\bibliography{references}

\end{document}
